% GENERATED FILE. Edit canonical lessons, then run npm run generate:books.
% canonical-source-hash: fnv1a64:f1511cc487de9c92
% canonical-lessons: RU-C405-shina, RU-C405-prokol, RU-C405-nasos, RU-C405-nakachivat, RU-C405-buksovat

\chapter{A flat tyre}
\label{ch:ru-a-flat-tyre}
\hlchaptermodality{405}

\begin{chapteropening}
\textbf{By the end of this chapter:} \emph{I can deal with a flat tyre: the tyre, the puncture, the pump, pumping it up and wheels that spin.}

Complete the last lesson of chapter 405: I can deal with a flat tyre: the tyre, the puncture, the pump, pumping it up and wheels that spin.
\end{chapteropening}

\section[\texorpdfstring{\ru{шина} (shína)}{шина (shína)}]{\ru{шина} (shína) — a tyre}

\label{lesson:RU-C405-shina}



\begin{quote}
Before the new one: say the Russian for to start, then the Russian for to stall.
\end{quote}

\subsection*{You'll want to know: \ru{шина}}
\textbf{\ru{шина}} (\emph{shína}) — \textquotedblleft{}a tyre\textquotedblright{}.

\begin{quote}
\ru{Передняя} \ru{шина} \ru{совсем} \ru{старая}.
\end{quote}

\begin{quote}
\emph{Perédnyaya shína sovsém stáraya.}
\end{quote}

\begin{quote}
The front tyre is quite old.
\end{quote}

Feminine. In everyday speech people also say \emph{\ru{колесо}} — a wheel.

\begin{cousinweb}
\emph{\ru{шина}} was borrowed from German \emph{Schiene} 'a rail, a strip' (EtymDB), from Old High German \emph{skina}; the word probably named the iron band round a cartwheel before the rubber tyre.
\end{cousinweb}

\begin{grammarlens}[title={What you've built}]
One more word for telling, explaining and describing at B1.
\end{grammarlens}

\subsection*{Your turn}
\begin{itemize}
\raggedright
  \item \emph{Say it:} \emph{shína}
  \item \emph{Say it:} \emph{shína}, once more
  \item \emph{Say it:} a sentence of your own with \emph{shína}
  \item \emph{Recall:} say the Russian for to start, then the Russian for to stall, then say \emph{shína} again
\end{itemize}

\begin{tcolorbox}[breakable,colback=teal!4,colframe=teal!35!black,title={Before you move on}]
What does \ru{шина} mean? (\textquotedblleft{}A tyre\textquotedblright{}.) Say it once more.
\end{tcolorbox}
\section[\texorpdfstring{\ru{прокол} (prokól)}{прокол (prokól)}]{\ru{прокол} (prokól) — a puncture}

\label{lesson:RU-C405-prokol}



\begin{quote}
Before the new one: say the Russian for to stall, then the Russian for a tyre.
\end{quote}

\subsection*{You'll want to know: \ru{прокол}}
\textbf{\ru{прокол}} (\emph{prokól}) — \textquotedblleft{}a puncture\textquotedblright{}.

\begin{quote}
\ru{Из}-\ru{за} \ru{прокола} \ru{мы} \ru{час} \ru{стояли} \ru{у} \ru{дороги}.
\end{quote}

\begin{quote}
\emph{Iz-za prokóla my chas stoyáli u dorógi.}
\end{quote}

\begin{quote}
Because of the puncture we stood by the road for an hour.
\end{quote}

Masculine. \emph{\ru{У} \ru{меня} \ru{прокол}} — I have a puncture.

\begin{cousinweb}
\emph{\ru{прокол}} is built on \emph{\ru{проколоть}} 'to pierce through', from \emph{\ru{про}-} 'through' and \emph{\ru{колоть}} 'to prick, to split'. EtymDB does not list \emph{\ru{прокол}}; its parts are ordinary Russian.
\end{cousinweb}

\begin{grammarlens}[title={What you've built}]
One more word for telling, explaining and describing at B1.
\end{grammarlens}

\subsection*{Your turn}
\begin{itemize}
\raggedright
  \item \emph{Say it:} \emph{prokól}
  \item \emph{Say it:} \emph{prokól}, once more
  \item \emph{Say it:} a sentence of your own with \emph{prokól}
  \item \emph{Recall:} say the Russian for to stall, then the Russian for a tyre, then say \emph{prokól} again
\end{itemize}

\begin{tcolorbox}[breakable,colback=teal!4,colframe=teal!35!black,title={Before you move on}]
What does \ru{прокол} mean? (\textquotedblleft{}A puncture\textquotedblright{}.) Say it once more.
\end{tcolorbox}
\section[\texorpdfstring{\ru{насос} (nasós)}{насос (nasós)}]{\ru{насос} (nasós) — a pump}

\label{lesson:RU-C405-nasos}



\begin{quote}
Before the new one: say the Russian for a tyre, then the Russian for a puncture.
\end{quote}

\subsection*{You'll want to know: \ru{насос}}
\textbf{\ru{насос}} (\emph{nasós}) — \textquotedblleft{}a pump\textquotedblright{}.

\begin{quote}
\ru{Хорошо}, \ru{что} \ru{у} \ru{нас} \ru{был} \ru{с} \ru{собой} \ru{насос}.
\end{quote}

\begin{quote}
\emph{Khoroshó, chto u nas byl s sobóy nasós.}
\end{quote}

\begin{quote}
It was a good thing we had a pump with us.
\end{quote}

Masculine. \emph{\ru{Велосипедный} \ru{насос}} — a bicycle pump.

\begin{cousinweb}
\emph{\ru{насос}} is built from \emph{\ru{на}-} and the root of \emph{\ru{сосать}} 'to suck': a pump sucks up air or water. EtymDB lists the word without a source.
\end{cousinweb}

\begin{grammarlens}[title={What you've built}]
One more word for telling, explaining and describing at B1.
\end{grammarlens}

\subsection*{Your turn}
\begin{itemize}
\raggedright
  \item \emph{Say it:} \emph{nasós}
  \item \emph{Say it:} \emph{nasós}, once more
  \item \emph{Say it:} a sentence of your own with \emph{nasós}
  \item \emph{Recall:} say the Russian for a tyre, then the Russian for a puncture, then say \emph{nasós} again
\end{itemize}

\begin{tcolorbox}[breakable,colback=teal!4,colframe=teal!35!black,title={Before you move on}]
What does \ru{насос} mean? (\textquotedblleft{}A pump\textquotedblright{}.) Say it once more.
\end{tcolorbox}
\section[\texorpdfstring{\ru{накачивать} (nakáchivat)}{накачивать (nakáchivat)}]{\ru{накачивать} (nakáchivat) — to pump up, to inflate}

\label{lesson:RU-C405-nakachivat}



\begin{quote}
Before the new one: say the Russian for a puncture, then the Russian for a pump.
\end{quote}

\subsection*{You'll want to know: \ru{накачивать}}
\textbf{\ru{накачивать}} (\emph{nakáchivat}) — \textquotedblleft{}to pump up, to inflate\textquotedblright{}.

\begin{quote}
\ru{Каждую} \ru{неделю} \ru{он} \ru{накачивает} \ru{шины} \ru{на} \ru{велосипеде}.
\end{quote}

\begin{quote}
\emph{Kázhduyu nedélyu on nakáchivayet shíny na velosipéde.}
\end{quote}

\begin{quote}
Every week he pumps up the tyres on his bicycle.
\end{quote}

Imperfective; perfective \emph{\ru{накачать}}.

\begin{cousinweb}
\emph{\ru{накачивать}} is built from \emph{\ru{на}-} and \emph{\ru{качать}} 'to swing, to pump'. EtymDB lists \emph{\ru{качать}} without a source, so its deeper origin is left open here.
\end{cousinweb}

\begin{grammarlens}[title={What you've built}]
One more word for telling, explaining and describing at B1.
\end{grammarlens}

\subsection*{Your turn}
\begin{itemize}
\raggedright
  \item \emph{Say it:} \emph{nakáchivat}
  \item \emph{Say it:} \emph{nakáchivat}, once more
  \item \emph{Say it:} a sentence of your own with \emph{nakáchivat}
  \item \emph{Recall:} say the Russian for a puncture, then the Russian for a pump, then say \emph{nakáchivat} again
\end{itemize}

\begin{tcolorbox}[breakable,colback=teal!4,colframe=teal!35!black,title={Before you move on}]
What does \ru{накачивать} mean? (\textquotedblleft{}To pump up, to inflate\textquotedblright{}.) Say it once more.
\end{tcolorbox}
\section[\texorpdfstring{\ru{буксовать} (buksovát)}{буксовать (buksovát)}]{\ru{буксовать} (buksovát) — to spin (of wheels), to get stuck}

\label{lesson:RU-C405-buksovat}



\begin{quote}
Before the new one: say the Russian for a pump, then the Russian for to pump up.
\end{quote}

\subsection*{You'll want to know: \ru{буксовать}}
\textbf{\ru{буксовать}} (\emph{buksovát}) — \textquotedblleft{}to spin (of wheels), to get stuck\textquotedblright{}.

\begin{quote}
\ru{В} \ru{снегу} \ru{колёса} \ru{долго} \ru{буксовали}.
\end{quote}

\begin{quote}
\emph{V snegú kolyósa dólgo buksováli.}
\end{quote}

\begin{quote}
In the snow the wheels spun for a long time.
\end{quote}

Imperfective; perfective \emph{\ru{забуксовать}}. Present \emph{\ru{буксует}}. Also figurative: \emph{\ru{проект} \ru{буксует}} — the project is stuck.

\begin{cousinweb}
The origin of \emph{\ru{буксовать}} is not settled. It is often linked with \emph{\ru{буксир}} 'a tug, a towrope', a sailors' word usually traced to Dutch; EtymDB does not list \emph{\ru{буксовать}}.
\end{cousinweb}

\begin{grammarlens}[title={What you've built}]
One more word for telling, explaining and describing at B1.
\end{grammarlens}

\subsection*{Your turn}
\begin{itemize}
\raggedright
  \item \emph{Say it:} \emph{buksovát}
  \item \emph{Say it:} \emph{buksovát}, once more
  \item \emph{Say it:} a sentence of your own with \emph{buksovát}
  \item \emph{Recall:} say the Russian for a pump, then the Russian for to pump up, then say \emph{buksovát} again
\end{itemize}

\begin{tcolorbox}[breakable,colback=teal!4,colframe=teal!35!black,title={Before you move on}]
What does \ru{буксовать} mean? (\textquotedblleft{}To spin (of wheels), to get stuck\textquotedblright{}.) Say it once more.
\end{tcolorbox}
